\chapter{Lições Aprendidas}

Esta seção tem como objetivo documentar as lições aprendidas ao longo do desenvolvimento do projeto do micromouse, de forma a consolidar os conhecimentos adquiridos, reconhecer os acertos da equipe e identificar os pontos passíveis de melhoria.

\section{Análise SWOT}

A fim de consolidar as lições aprendidas ao longo do desenvolvimento do micromouse, realizou-se uma análise SWOT do produto, contemplando os fatores internos (pontos fortes e fracos da equipe e do projeto) e os fatores externos (oportunidades e ameaças relacionadas ao contexto do desenvolvimento). A Tabela~\ref{tab:swot} sintetiza os principais pontos identificados pela equipe.

\begin{table}[H]
\centering
\footnotesize
\caption{Análise SWOT do produto.}
\begin{tabular}{|c|c|c|} \hline
 & \textbf{Pontos fortes} & \textbf{Pontos fracos}  \\ \hline 
{\parbox{0.1\textwidth}{ \textbf{Fatores \\ internos}}}  &
{\parbox{0.4\textwidth}{ Forças: 
\begin{itemize}
    \item Chassi leve, com bom desempenho geral
    \item Bateria eficiente e boa autonomia
    \item Lógica planejada previamente, agilizando diagnóstico de erros
    \item Inerciômetro como redundância, mitigando falha de encoder
    \item Experiência prévia da equipe em tecnologias específicas
    \item Arquitetura em React.js modular e escalável
    \item Comunicação em tempo real via WebSocket
    \item Extras de valor agregado, como carrinho em 3D
    \item Agilidade no dimensionamento da alimentação elétrica
\end{itemize}}}
&
{\parbox{0.4\textwidth}{ Fraquezas:
\begin{itemize}
    \item Largura/formato do chassi maior que o ideal, dificultando manobras nas curvas
    \item Queima de componentes durante o processo
    \item Aumento do orçamento previsto devido à reposição dos componentes danificados
    \item Eletrônica mais sensível devido ao uso de jumpers em placa perfurada
    \item Banco de dados dependente de conexão à internet, exigindo solução de backup offline
    \item Predominância de testes isolados em detrimento de testes de integração, gerando imprevistos na integração final
\end{itemize}}}
\\ \hline

{\parbox{0.1\textwidth}{ \textbf{Fatores \\ externos}}}  &
{\parbox{0.4\textwidth}{ Oportunidades:
\begin{itemize}
    \item Reaproveitamento do chassi, algoritmo de navegação e sistema de energia como base para outros projetos acadêmicos, TCCs ou iniciação científica
    \item Adaptação dos módulos desenvolvidos (gerenciamento de energia, navegação) para outros robôs móveis, como VANTs ou rovers
\end{itemize}}}
&
{\parbox{0.4\textwidth}{ Ameaças:
\begin{itemize}
    \item Prazo apertado devido a imprevistos técnicos
    \item Falhas pequenas acumuladas dificultando diagnóstico na integração
    \item Alto escopo de entregas em curto prazo
    \item Dificuldade em organizar o GitHub em relação à EAP
    \item Demora na integração com hardware, atrasando a telemetria
\end{itemize}}}
\\ \hline
\end{tabular}
\label{tab:swot}
\end{table}


\section{Planejado x realizado}

\subsection{Prazos}

\subsection{Orçamento}

No início do projeto, foi pensado um orçamento máximo de 1.000 reais. A intenção era trabalhar com um valor abaixo do limite, porém, devido ao caráter iterativo do projeto, surgiram custos adicionais, já que diversas soluções precisaram ser testadas e ajustadas até chegar ao resultado desejado. Por outro lado, no início do desenvolvimento, a equipe conseguiu algumas peças emprestadas, o que proporcionou uma maior folga orçamentária. Dessa forma, mesmo com os imprevistos, o valor total permaneceu dentro do limite estabelecido, fechando em 853,88 reais, sendo que a maior parte do investimento foi direcionada aos componentes eletrônicos.

\subsection{Escopo}

No que diz respeito ao escopo do projeto, os objetivos definidos foram majoritariamente cumpridos e o micromouse foi construído conforme planejado, 
integrando todas as áreas devidamente. Porém, alguns componentes e funcionalidades não foram incorporados à versão final do projeto, na área de hardware, por exemplo,
o giroscópio MPU-6050 foi removido devido ao ruído eletromagnético que comprometeu a leitura dos valores. Além disso, a capacidade de medição de direção 
da rotação dos motores foi comprometida e o fio C2 dos encoders não foi utilizado. Em firmware, os switches de seleção e o botão de liga/desliga não foram integrados ao sistema devido
ao espaço limitado da placa.


\section{Projeto}

A partir da experiência acumulada durante as etapas de desenvolvimento, integração e testes do micromouse, foram identificados os principais pontos fortes e fracos do projeto, descritos a seguir.

\subsection{Pontos Fortes}

\begin{enumerate}
    \item Chassi leve e de montagem simples, facilitando ajustes e manutenção ao longo do projeto.
    \item Uso do inerciômetro como redundância, essencial diante da falha em um dos encoders.
    \item Arquitetura em React.js e comunicação via WebSocket garantiram um frontend modular e responsivo.
    \item Experiência prévia da equipe em tecnologias específicas acelerou o desenvolvimento geral.
\end{enumerate}

\subsection{Pontos fracos}

\begin{enumerate}
    \item Formato do chassi dificultou manobras nas curvas do labirinto.
    \item Queima de componentes durante o processo, elevando o orçamento previsto do projeto.
    \item Predomínio de testes isolados sobre testes de integração, gerando imprevistos na integração final.
    \item Engajamento desigual entre os membros da equipe, dificultando o cumprimento dos prazos.
\end{enumerate}

\section{Recomendações para projetos futuros}

\begin{enumerate}
    \item Priorizar entregas críticas do caminho principal antes de 
    funcionalidades secundárias, evitando acúmulo de pendências próximo 
    ao prazo final.
    
    \item Melhorar a comunicação entre as subequipes, estabelecendo 
    reuniões regulares de alinhamento para reduzir retrabalho causado 
    por decisões tomadas de forma isolada.
    
    \item Concluir cada entrega de forma sólida antes de avançar para 
    a próxima, evitando retrabalho comprometer a integração final do sistema.
    
    \item Adotar revisão em dupla para todas as entregas técnicas, 
    reduzindo erros simples e aumentando a qualidade do código e da 
    documentação produzidos.
\end{enumerate}

\section{Desempenho dos fornecedores}


Esta seção apresenta uma avaliação dos fornecedores utilizados ao longo do desenvolvimento do projeto, classificados conforme o desempenho apresentado em relação ao esperado pela equipe.

\subsection{Fornecedores com desempenho acima do esperado}

Fornecedores que a equipe certamente convidaria para um próximo projeto.

\begin{itemize}
    \item \textbf{HU Infinito} -- Ampla variedade de produtos, atendimento de qualidade e entregas sempre pontuais.
    \item \textbf{SIA Parafusos e Ferramentas} -- Atendimento rápido, apesar de nem todos os modelos necessários estarem disponíveis em estoque.
\end{itemize}

\subsection{Fornecedores com desempenho conforme esperado}

Fornecedores que a equipe provavelmente convidaria para um próximo projeto.

\begin{itemize}
    \item \textbf{RoboCore} -- Pedido entregue com boa qualidade, porém com atraso em relação ao prazo combinado.
\end{itemize}

\subsection{Fornecedores com desempenho abaixo do esperado}

Fornecedores que a equipe não convidaria para um próximo projeto.

\begin{itemize}
    \item \textbf{Casa dos Parafusos} -- Variedade adequada de produtos, porém atendimento demorado.
    \item \textbf{Dr. Eletrônico} -- Entrega não prática e atendimento confuso durante o processo de compra.
\end{itemize}


% Resolvi não pagar os comentarios abaixo por enquanto para fins de consulta


% \chapter{Lições Aprendidas}

% \textcolor{red}{Documentar as lições aprendidas de modo a aperfeiçoar os processos e evitar que os erros e problemas encontrados se repitam em futuros projetos.}

% \section{Análise SWOT}

% \textcolor{red}{A análise SWOT (\textit{Strengths}, \textit{Weaknesses}, \textit{Opportunities}, \textit{Threats}) é usada para identificar os pontos fortes e fracos do produto, e as principais oportunidades e ameaças normalmente associadas aos pontos fortes e fracos identificados. Preencha a Tabela \ref{tab:swot} de forma correspondente.}

% \begin{table}[ht]
% % \label{tab:swot}
% \centering
% \caption{Análise SWOT do produto.}
% \begin{tabular}{|c|c|c|} \hline
%  & \textbf{Pontos fortes} & \textbf{Pontos fracos}  \\ \hline 
%     {\parbox{0.1\textwidth}{ \textbf{Fatores \\ internos}}}  &
%     {\parbox{0.4\textwidth}{ Forças: 
%         \\ \textcolor{red}{Descreva seus diferenciais competitivos: aquilo que faz de você melhor que seus concorrentes.}
%         \begin{itemize}
%             \item Lorem ipsum
%             \item Lorem ipsum
%         \end{itemize} }} & 
%     {\parbox{0.4\textwidth}{ Fraquezas:
%         \\ \textcolor{red}{Descreva suas principais deficiências: aquilo que faz de você pior que seus concorrentes, e pontos a serem melhorados ou revistos.}
%         \begin{itemize}
%             \item Lorem ipsum
%             \item Lorem ipsum
%         \end{itemize} }}\\ \hline
%     {\parbox{0.1\textwidth}{ \textbf{Fatores \\ externos}}}  &
%     {\parbox{0.4\textwidth}{ Oportunidades:
%         \\ \textcolor{red}{Descreva as principais oportunidades identificadas: eventos potenciais que podem gerar grandes benefícios.}
%         \begin{itemize}
%             \item Lorem ipsum 
%             \item Lorem ipsum
%         \end{itemize} }} & 
%     {\parbox{0.4\textwidth}{ Ameaças:
%         \\ \textcolor{red}{Descreva as principais ameaças identificadas: eventos potenciais que podem causar um grande estrago ao seu negócio.}
%         \begin{itemize}
%             \item Lorem ipsum 
%             \item Lorem ipsum
%         \end{itemize} }}\\ \hline
% \end{tabular}
% \label{tab:swot}
% \end{table}


% \section{Planejado x realizado}
% \subsection{Objetivos}
% \textcolor{red}{Os objetivos foram alcançados? Responda as questões e comente os pontos mais relevantes – até 700 caracteres, considerando os espaços.}
% \subsection{Prazos}
% \textcolor{red}{O projeto foi entregue dentro do prazo? Responda as questões e comente os pontos mais relevantes – até 700 caracteres, considerando os espaços.}
% \subsection{Orçamento}
% \textcolor{red}{O projeto foi entregue dentro do orçamento? Responda as questões e comente os pontos mais relevantes – até 700 caracteres, considerando os espaços.}
% \subsection{Escopo}
% \textcolor{red}{O  projeto atendeu ao escopo? Responda as questões e comente os pontos mais relevantes – até 700 caracteres, considerando os espaços.}


% \section{Projeto}
% \textcolor{red}{Comente os pontos mais relevantes a serem aperfeiçoados ou adotados em próximos projetos.}
% \subsection{Pontos fortes}
% \begin{enumerate}
% 	\item \textcolor{red}{Até 200 caracteres, considerando os espaços.}
% 	\item \textcolor{red}{Até 200 caracteres, considerando os espaços.}
% 	\item \textcolor{red}{Até 200 caracteres, considerando os espaços.}
% 	\item \textcolor{red}{Até 200 caracteres, considerando os espaços.}
% \end{enumerate}
% \subsection{Pontos fracos}
% \begin{enumerate}
% 	\item \textcolor{red}{Até 200 caracteres, considerando os espaços.}
% 	\item \textcolor{red}{Até 200 caracteres, considerando os espaços.}
% 	\item \textcolor{red}{Até 200 caracteres, considerando os espaços.}
% 	\item \textcolor{red}{Até 200 caracteres, considerando os espaços.}
% \end{enumerate}

% \section{Recomendações para projetos futuros}

% \begin{enumerate}
% 	\item \textcolor{red}{Até 200 caracteres, considerando os espaços.}
% 	\item \textcolor{red}{Até 200 caracteres, considerando os espaços.}
% 	\item \textcolor{red}{Até 200 caracteres, considerando os espaços.}
% 	\item \textcolor{red}{Até 200 caracteres, considerando os espaços.}
% \end{enumerate}

% \section{Questões em aberto}

% \textcolor{red}{Usar caso haja alguma questão pendente em relação às entregas do projeto (Ex.: Requisitos não entregues, melhoria na programação do software, troca de sensores, entre outros.}

% \begin{enumerate}
% 	\item \textcolor{red}{Até 200 caracteres, considerando os espaços.}
% 	\item \textcolor{red}{Até 200 caracteres, considerando os espaços.}
% 	\item \textcolor{red}{Até 200 caracteres, considerando os espaços.}
% 	\item \textcolor{red}{Até 200 caracteres, considerando os espaços.}
% \end{enumerate}

% \section{Desempenho dos fornecedores}

% \subsection{Fornecedores com desempenho acima do esperado}
% \textcolor{red}{Fornecedores que você certamente convidaria para um próximo projeto. Citar AAAA – justificativa. Até 200 caracteres, considerando os espaços.}

% \subsection{Fornecedores com desempenho conforme esperado}
% \textcolor{red}{Fornecedores que você provavelmente convidaria para um próximo projeto. Citar AAAA – justificativa. Até 200 caracteres, considerando os espaços.}

% \subsection{Fornecedores com desempenho abaixo do esperado}
% \textcolor{red}{Fornecedores que você não convidaria para um próximo projeto. Citar AAAA – justificativa. Até 200 caracteres, considerando os espaços.}